\documentclass[11pt]{article}

\usepackage[margin=0.95in]{geometry}
\usepackage[T1]{fontenc}
\usepackage{lmodern}
\usepackage{microtype}
\usepackage{amsmath,amssymb}
\usepackage{xcolor}
\usepackage{enumitem}
\usepackage{fancyhdr}
\usepackage{lastpage}
\usepackage{graphicx}
\usepackage{booktabs}
\usepackage{array}
\usepackage{tabularx}
\usepackage[colorlinks=true,linkcolor=Primary,citecolor=Primary,urlcolor=Secondary]{hyperref}
\usepackage[most]{tcolorbox}

\definecolor{Primary}{HTML}{2C5F52}
\definecolor{Secondary}{HTML}{B8860B}
\definecolor{ProblemBg}{HTML}{FAF7F0}
\definecolor{TitleBg}{HTML}{1F4037}

\newcommand{\Course}{COMPSCI 566}
\newcommand{\Assignment}{Project Midterm Report}
\newcommand{\ProjectTitle}{Monocular Depth Estimation: Benchmarking and a Defocus-Informed Refinement}
\newcommand{\Date}{\today}
\newcommand{\captext}[1]{{\fontsize{8.5}{10}\selectfont\itshape #1\par}}

\setlength{\parindent}{0pt}
\setlength{\parskip}{2pt}
\linespread{1.0}
\setlist[enumerate]{label=\textbf{\arabic*.}, leftmargin=*, itemsep=0.25em, topsep=0.25em}
\setlist[itemize]{leftmargin=*, itemsep=0.1em, topsep=0.1em}

\pagestyle{fancy}
\fancyhf{}
\lhead{\textcolor{Primary}{Chrysafis / Zhou}}
\chead{\textcolor{Primary}{\Course}}
\rhead{\textcolor{Primary}{\Assignment}}
\cfoot{\textcolor{Primary!85!black}{Page \thepage\ of \pageref{LastPage}}}
\setlength{\headheight}{24pt}
\renewcommand{\headrulewidth}{0.7pt}
\renewcommand{\headrule}{\hbox to\headwidth{\color{Primary!55}\leaders\hrule height \headrulewidth\hfill}}

\tcbset{
  hwbox/.style={
    enhanced,
    breakable,
    colback=ProblemBg,
    colframe=Primary!85!black,
    boxrule=0.85pt,
    arc=2mm,
    left=2.6mm,
    right=2.6mm,
    top=1.3mm,
    bottom=1.3mm,
    before skip=2pt,
    after skip=2pt,
    fonttitle=\bfseries,
    coltitle=Primary!95!black,
    attach boxed title to top left={xshift=2mm,yshift*=-2mm},
    boxed title style={
      colback=Secondary!18,
      colframe=Primary!80!black,
      boxrule=0.8pt,
      arc=1mm,
      left=1.5mm,
      right=1.5mm,
      top=0.8mm,
      bottom=0.8mm,
    },
  },
}
\newtcolorbox{reportbox}[2][]{hwbox, title={#2}, #1}

\begin{document}

\begin{tcolorbox}[
  enhanced,
  arc=2.5mm,
  boxrule=0pt,
  colback=TitleBg,
  colframe=TitleBg,
  borderline south={1.5pt}{0pt}{Secondary},
  title={\textcolor{white}{\Assignment\ \textbullet\ \Date}},
  coltitle=white,
  fonttitle=\bfseries,
  before skip=0pt,
  after skip=5pt
]
  {\color{white}
  \begin{tabular*}{\textwidth}{@{\extracolsep{\fill}}ll}
    \textbf{Fall 2026} & \textbf{Course:} \Course
  \end{tabular*}\\[3pt]
  {\color{white}\textbf{\ProjectTitle}}\\[2pt]
  {\color{white}\small Dimitri Chrysafis (chrysafis) \quad Vincent Zhou (vzhou2)}
  }
\end{tcolorbox}

% Required by the syllabus: brief summary of current progress, current
% results, difficulties during implementation, and how the proposal has
% changed in light of progress.

\begin{reportbox}{Progress Summary}
The full pipeline is implemented and evaluated end to end. Phase 1 is
complete: DA-V2 ViT-S is benchmarked on all 654 NYU-Depth v2 test
images, and the classical depth-from-focus baseline (modified
Laplacian focus measure, parabolic sub-frame interpolation, confidence
map) runs on rendered focal stacks and all 120 DDFF test stacks.
Phase 2 is implemented and in ablation: the robust scale fit and the
regularized local correction are working on the synthetic scenes, and
the combined system runs qualitatively on real DDFF stacks.
\end{reportbox}

\begin{reportbox}{Current Results}
NYU-Depth v2 test split (654 images), DA-V2 ViT-S, two alignment
variants against ground truth:
\begin{center}
\begin{tabular}{lccc}
\toprule
Alignment & AbsRel & RMSE (m) & $\delta_1$ \\
\midrule
Least-squares, scale+shift & 0.095 & 0.431 & 0.912 \\
Median ratio, scale only & 0.507 & 2.376 & 0.376 \\
\bottomrule
\end{tabular}
\end{center}
The affine fit matters: dropping the shift term roughly doubles AbsRel,
which supports modeling $q = s\hat d + t$ rather than rescaling alone.

Synthetic focal stacks (10 frames, focus 0.9--7\,m, thin-lens blur
model, means over four scenes), metric depth scored without alignment:
\begin{center}
\begin{tabular}{lccc}
\toprule
Method & AbsRel & RMSE (m) & $\delta_1$ \\
\midrule
Focus baseline (DFF) & 0.202 & 1.521 & 0.706 \\
DA-V2 aligned (oracle) & 0.255 & 1.296 & 0.567 \\
DA-V2 + scale fit from focus & 0.284 & 1.508 & 0.502 \\
Full refinement & 0.248 & 1.530 & 0.560 \\
\bottomrule
\end{tabular}
\end{center}
The focus baseline is accurate where texture exists
($\delta_1$ 0.96--1.00 on textured regions) but degrades on flat
regions, exactly where the confidence map reads low
($\bar w \approx 0.06$--0.25 textureless vs.\ 0.26--0.66 textured).
On the calibrated scenes' textured pixels the refinement improves on
the focus baseline (AbsRel 0.154 vs.\ 0.158) and beats the affine-only
fit overall (0.248 vs.\ 0.284), i.e.\ the local correction $r$
contributes beyond the global scale terms. DA-V2 is weaker than the focus
baseline on our synthetic scenes even with oracle alignment ---
plausible, since the renderer is far from its training distribution ---
which makes this benchmark a useful stress test rather than a real-data
statement.
\end{reportbox}

\begin{reportbox}{Difficulties}
\begin{itemize}
\item \textbf{DDFF labels.} The public test archive ships 120 focal
  stacks but no ground truth (withheld for the benchmark server), and
  per-frame focus distances are not recorded. We reconstructed focus
  distances from the refocus geometry in the paper (disparities linear
  in $[0.02,0.28]$~px, $z \approx 0.14/d$) and limited DDFF to
  qualitative evaluation; measured confidences on the real Lytro stacks
  are low ($\bar w \approx 0.03$--0.05), and scale calibration was
  rejected on half of the stacks we tried.
\item \textbf{NYU test split sourcing.} Both canonical sources were
  unreachable (the DenseDepth mirror returns 403 and the NYU host
  refuses connections); we used a mirror of the same 654-image Eigen
  test split and verified image count, units, and depth range.
\item \textbf{Defocus veiling.} In the layered renderer, blurred
  near-field texture bleeds over flat surfaces, which lifts the focus
  measure there (the synthetic corridor's flat ceiling reaches
  $w \approx 0.4$). This is a genuine failure mode of the focus
  measurement rather than a renderer artifact and it caps how far we
  can trust $w$ at occlusion boundaries.
\item \textbf{Rejected fits.} On scenes where the learned prediction
  anticorrelates with the focus estimate, the scale fit is degenerate
  ($s \le 0$); we fall back to the bounded uncalibrated output per the
  proposal rather than emit bad metric depth.
\end{itemize}
\end{reportbox}

\begin{reportbox}{Changes from Proposal}
DDFF 12-Scene moved from a quantitative to a qualitative dataset: the
verification the proposal made conditional showed the released test
stacks carry no labels. Rendered scenes now carry the quantitative
refinement evaluation, which also gives us controlled ablations. The
method itself is unchanged; we added a degenerate-fit guard so
uncalibrated output is reported as relative depth.
\end{reportbox}

\begin{reportbox}{Next Steps}
\begin{itemize}
\item First-pass ablations are in (stack size, confidence weighting,
  $\lambda$, $\mu$); next is a tuning/held-out scene split so the
  reported numbers are not tuned on the evaluation scenes.
\item Extend the DDFF qualitative evaluation and, time permitting,
  phone-captured focal stacks as a second real source.
\item Finish the results figures and tables for the project webpage;
  keep the presentation deck in step with the final numbers.
\end{itemize}
\end{reportbox}

\end{document}
